% ---------------------------------------------------------------------------
% Answers the four questions, separates what is established from what is
% assumed, and says what the next measurement would be.
% ---------------------------------------------------------------------------

\section{Conclusion}
\label{sec:conclusion}

\paragraph{The four questions.}
\textbf{RQ1.} Both languages admit deterministic one-token-lookahead parsers, with an
asymmetry that is itself a finding. The inferred Datadog grammar --- the first published
grammar for that language --- is conflict-free. PromQL has $2$ shift/reduce conflicts,
both on the parenthesis that begins either a vector-matching label list or a parenthesised
operand, and both forced by that list's optionality rather than by how we layered the
rules; so RQ1's answer for PromQL is ``LALR(1) modulo one documented resolution'', for
this grammar, with no claim that none does better. Neither answer is where the switching
cost lives, which is itself worth knowing.

\textbf{RQ2.} The largest exactly-translating fragment we could build covers $96.3\%$ of
Datadog query occurrences, and the soundness theorems of \S\ref{sec:kernel-soundness} and
\S\ref{sec:binop-soundness} between them cover every expression the translator emits, so
nothing is shipped that is not proved. Outside the fragment, $1.2\%$ of the corpus is
beyond reach --- four constructs, two of them information barriers rather than
expressiveness gaps --- against $2.4\%$ that is work this translator has not done, a
ceiling of $98.7\%$. In the reverse direction, on the sub-problem of tag filtering, the
obstructions are theorems: a glob's language is a single word or unbounded in length, and
three impossibility results follow. With a stated bound on union size they account for
every refusal we measure --- $37.3\%$ of the matcher occurrences whose pattern is
concrete --- and $10$ of those need no bound at all.

\textbf{RQ3.} Coverage on published dashboards is the $96.3\%$ above, and the shape of
the shortfall matters more than its size: coverage against implementation effort is a
staircase rather than a curve, because a query translates only when every feature it uses
is supported, while the popularity that would guide an implementation order is marginal.
Half the work does not port half the estate.

\textbf{RQ4.} No. A dashboard query does not determine its own meaning, and this is the
answer we would now put first. Of the translations we emit, $96.9\%$ depend on a rollup
interval the source query leaves to the dashboard's time range, and $100\%$ on a metric
type no query in either language declares. Neither is a defect of one vendor: $88.6\%$ of
Datadog query occurrences and $82.5\%$ of PromQL ones carry at least one template
variable, so in both languages the object a migration ports is a schema plus a
resolution layer that belongs to the dashboard tool. One of the two assumptions is
checkable, because the repository the dashboards come from declares each metric's type,
and where it can be checked our inference is wrong for $6.5\%$ --- $491$ occurrences,
more than the whole inexpressible column. $\kappa$ counts those as translations.
Coverage and correctness differ, by a measured amount, and no coverage number could have
told us: a query that does not fix its own meaning parses and translates like any
other.

\noindent
What each of those answers rests on --- what is proved, what is verified against a real
server, what is measured against the vendor, and what is still assumed --- is set out
\apref{app:experiments}. Three of the eight vendor measurements contradicted what this
paper had assumed before asking.

\iffull% The inventory of evidence that closed the conclusion: proved, verified against a
% real server, read and written down, measured against the vendor, and assumed.
\subsection{What rests on what}
\label{app:evidence}

Three things in this paper are proved: the directional information barrier, whose
reverse half is Proposition~\ref{prop:reverse} and holds for every possible translator
and not only ours; the exactness of the rollup translations on a
stated class, with the preconditions enumerated as the hypotheses of a theorem; and the
non-expressibility results of \S\ref{sec:regular}. Three things are verified against a
real Prometheus~3.15.0 rather than a model, with the responses cached and
regression-tested: the PromQL semantics the rollup rules depend on, the anchoring of
label-matcher regular expressions, and the series pairing that
Theorem~\ref{thm:arith} rests on. Two further readings were taken from a real server
and written down as prose rather than cached --- the parse error that decides
\S\ref{sec:grammars}'s two conflicts, and the rejection that
Remark~\ref{rem:anchored-flag} turns on --- so they are evidence of a weaker kind than
the three above, and the difference is in the artefact rather than in the reading.
Nothing is left that was read from source and not then measured: Datadog's reset policy was taken from its
open-source agent at a pinned commit and afterwards put to a running one.

Eight things are measured against the vendor itself rather than assumed: the bucket
convention (left-closed and named by its start, the opposite of what this work assumed
until it asked), the automatic rollup interval as a function of the query range, the
default time aggregator of a bare query on a counter, whether a glob's \texttt{*} matches
the empty string, where \texttt{.as\_rate()} sits relative to the rollup and to the space
aggregation, what \texttt{.weighted()} computes, what \texttt{top(...)} ranks by and over
which window, and the reset policy --- read first from
the agent's source and then put to a running agent, which agreed with the reading. Five
of the eight confirmed what the paper had assumed. Three did not: the bucket convention;
\texttt{.as\_rate()}, which turned out to divide the aggregated result rather than
each series, and whose presence in a query turned out not to establish the metric's type
at all, because Datadog accepts the modifier on a gauge and ignores it; and
\texttt{.weighted()}, carried here as the largest single entry in the inexpressible
column until the account was asked and returned the same number with the modifier and without it. The last of the
eight decides which column $139$ corpus occurrences belong in, and is the only one
measured on both sides: over one workload, the candidate expression against a real
Prometheus in $20$ cases and \texttt{top()} itself against Datadog in $22$, each
agreeing with the documented rule throughout.

One assumption is settled from the vendor's source tree rather than from its API, and
it is the one that held of all the output. Every translation carries
\texttt{METRIC\_TYPE\_UNKNOWN} because a dashboard does not record whether a metric is
a counter or a gauge --- but \texttt{integrations-core} declares the type of each
metric it ships, in a file beside the dashboards that the corpus fetch had passed over.
For the $89.5\%$ of translated occurrences whose metric is declared, the front end's
inference is wrong on $6.5\%$: $491$ occurrences against the $0.8\%$ this
paper calls beyond reach (Remark~\ref{rem:type-measured}). $\kappa$ counts them as
translations, so the gap between coverage and correctness has a number for the first
time, and it is not zero.

One thing that was assumed throughout is now measured. The counter rules are checked
against Datadog itself, by replaying a real request trace into both systems and comparing
bucket by bucket (\S\ref{sec:replay}): both exact forms are exact and \texttt{rate} is
not. The one part a replay cannot reach is the reset policy, because a request counter
does not restart in half an hour; that took a counter built to restart, scraped across
the reset by a real agent and a real Prometheus at once, and it is the last assumption of
\S\ref{sec:ingestion} to be retired (Remark~\ref{rem:rho}).

The same construction has since been run at corpus scale, and it is the reason $\kappa$
can now be read as a statement about translations that were checked rather than merely
emitted. Abstracting the identifiers away --- and nothing a query would emit differently ---
collapses the corpus to $332$
distinct query shapes, few enough that a single trial account holds every series they
need; each is instantiated, written to both systems and compared bucket by bucket
(\S\ref{sec:differential}). It found four translation defects that the unit tests and the
single-trace replay had both missed --- and, six times over, found its own harness
first, which is the more useful lesson about differential rigs than the agreement rate
is.
\fi

\paragraph{What we would do next.}
The reverse translator: \S\ref{sec:regular} suggests the interesting results are there, and a
$\kappa$ for that direction would turn this paper's asymmetry claim from an argument into
a measurement. Further out, the method is language-agnostic wherever a language denotes a
function of a database at an instant, which brings LogQL within reach and leaves SPL
outside it.

\paragraph{The answer.}
Is the switching cost a language-design artefact? Partly, and not where we expected. In
the direction teams travel it is mostly engineering --- a finding this paper four times
nearly got wrong in the more quotable direction, through a structural error in the
translator, a rule it had not implemented, a modifier it never asked the vendor about,
and a classification that contradicted its own text. The artefact in that direction is
subtler than an expressiveness gap and is not a property of either language on its own: a
dashboard query does not say what it means, so every translation we emit is exact only
relative to context that arrives from the layer above both languages. The expressiveness
gap is real, provable, and in the direction teams do not travel.

That is the asymmetry --- leaving costs work, returning costs something work cannot buy
--- and it is worth being exact about whom it binds, because it is not the mechanism the
lock-in literature usually describes. We measured the exit and found it cheap. What is
expensive is the return, so nobody is held by what we measured; a team that wants to
leave, can. The cost falls instead on the team that has left and wants back, and on the
team that will not leave because it cannot price the option to return. That is a cost of
irreversibility rather than of captivity, and it is paid by a different decision at a
different time. Our two directions are also drawn from different populations ---
$\kappa$ over whole queries against a share of concrete tag-filter patterns, with no
reverse translator built --- so the asymmetry is an argument from two measurements and
not yet one measurement of two directions. Making it the latter is the next thing to do.
